% AI-assisted manuscript source; responsible author: Carptopus.
\documentclass[11pt]{article}
\usepackage[a4paper,margin=30mm]{geometry}
\usepackage[T1]{fontenc}
\usepackage{lmodern}
\usepackage{microtype}
\usepackage{amsmath,amssymb,amsthm,mathtools}
\usepackage{enumitem}
\usepackage{fvextra}
\usepackage{xcolor}
\usepackage[colorlinks=true,linkcolor=blue!55!black,citecolor=blue!55!black,urlcolor=blue!65!black]{hyperref}
\usepackage{fancyhdr}

\newtheorem{theorem}{Theorem}[section]
\newtheorem{lemma}[theorem]{Lemma}
\newtheorem{corollary}[theorem]{Corollary}
\theoremstyle{remark}
\newtheorem{remark}[theorem]{Remark}

\pagestyle{fancy}
\fancyhf{}
\fancyhead[L]{Carptopus}
\fancyhead[R]{Spectral separation and the nSSP}
\fancyfoot[C]{\thepage}
\setlength{\headheight}{14pt}
\setlength{\parindent}{0pt}
\setlength{\parskip}{0.55em}
\setlength{\emergencystretch}{2em}
\setlist{nosep,leftmargin=*}

\title{Spectral separation and the non-symmetric strong spectral property\\for looped double paths}
\author{Carptopus\\\href{mailto:carptopus@163.com}{carptopus@163.com}}
\date{19 August 2026}

\hypersetup{
  pdftitle={Spectral separation and the non-symmetric strong spectral property for looped double paths},
  pdfauthor={Carptopus},
  pdfsubject={Non-symmetric strong spectral property for looped double paths},
  pdfkeywords={nSSP, inverse eigenvalue problem, looped double path, Jacobi matrix, spectral separation}
}

\begin{document}
\maketitle

\begin{center}
\small\itshape
Public Beta v0.1. The theorem chain has passed an internal writing-admission audit,
exact computational calibration, and a manuscript-level zero-trust audit. External mathematical
review and journal peer review are pending. Priority statements are qualified by
``to the best of our knowledge.''
\end{center}
\vspace{0.8em}

\begin{abstract}


Let $P_{n,L}$ be the double path on $[n]$, with a loop at precisely the vertices in
$L\subseteq[n]$. We classify when this zero--nonzero pattern allows the non-symmetric strong spectral
property (nSSP). The main structural step concerns a real symmetric irreducible tridiagonal matrix with
exactly one nonzero diagonal entry. We prove that it has the nSSP if and only if the two zero-diagonal
Jacobi arms obtained by deleting the looped vertex have disjoint spectra. The forward obstruction is an
explicit commuting spectral-projector certificate; the reverse implication follows from parity separation
of the characteristic polynomial and a rooted-continuant recovery argument. Scaling one arm avoids the
finitely many nonzero spectral collisions whenever parity does not force a common zero eigenvalue. Using
the Superpattern Lemma and the known nonexistence theorem for odd double paths whose loops lie only on
even vertices, we obtain

\[
  P_{n,L}\text{ allows the nSSP}
  \quad\Longleftrightarrow\quad
  L\neq\varnothing\ \text{ and }\bigl(n\text{ is even or }L\text{ contains an odd vertex}\bigr).
\]

In particular, the answer to Question 5.11 of Koljan\v{c}i\'{c} and Oblak is affirmative. The result classifies
allowance, not the finer distinction between patterns that require the nSSP and those that merely allow it.

\end{abstract}

\section{Introduction}


For a real matrix $A$, the non-symmetric strong spectral property requires that the simultaneous
linear equations

\[
  A\circ X=0,
  \qquad
  [A,X^{\mathsf T}]=0
\]

have only the zero solution. Here $\circ$ is the Hadamard product and
$[A,B]=AB-BA$. This property is a transversality condition for inverse eigenvalue problems and is
preserved, in an appropriate existence sense, when nonzero entries are added to a pattern through the
Superpattern Lemma [3, 4].

Koljan\v{c}i\'{c} and Oblak [4] study the nSSP for \textbf{double paths}: irreducible tridiagonal zero--nonzero
patterns with a prescribed set of nonzero diagonal positions. They show that loop placement creates all
three behaviours relevant to their classification problem: a pattern may require the nSSP, may allow but
not require it, or may fail to allow it. They ask in Question 5.11 whether every even-order double path
with at least one loop allows the nSSP.

We answer that question and, more generally, classify the allow/not-allow boundary for every order and
every loop assignment. The key is to use a symmetric matrix inside the generally non-symmetric pattern.
For a matrix with a single loop, deleting the looped vertex leaves two zero-diagonal Jacobi-type arms.
Their spectra provide an exact certificate:

\begin{quote}
\emph{the symmetric single-loop matrix has the nSSP exactly when the two arm spectra are disjoint.}
\end{quote}

This statement is closely related to the classical three-spectra inverse problem for Jacobi matrices.
Michor and Teschl [5] showed that the spectrum of a Jacobi matrix together with the spectra of the two
submatrices obtained by deleting an interior vertex gives unique reconstruction without additional mixing
data precisely when the two submatrix spectra are disjoint. Spectral separation and continuant recovery
are therefore not new by themselves. The contribution here is their connection to the nSSP pattern
functionals, an explicit nSSP obstruction when the arm spectra meet, and the resulting parity
classification of loop assignments.

The proof has three conceptual parts.

\begin{enumerate}
\item Simple spectrum identifies every commuting matrix with a polynomial in $A$. The nSSP equations
   then become exactly $n$ linear functionals on $\mathbb R[A]$, equivalently the Jacobian of the
   characteristic coefficients with respect to the one loop value and the $n-1$ edge weights.
\item The zero-diagonal path polynomial and the deleted-vertex polynomial have opposite parity. An
   infinitesimal isospectral perturbation therefore separates into two equations. If the arm polynomials
   are coprime, rooted continuants recover every edge-weight square to first order. If the arm spectra
   meet, spectral projectors give a nonzero commuting matrix supported outside the pattern.
\item Parity forces a common zero eigenvalue exactly when $n$ is odd and the unique loop is at an even
   vertex. In all other cases, scaling one arm avoids the finitely many remaining collisions.
\end{enumerate}
Our final theorem uses two prior results exactly at their established scope: loopless patterns do not
allow the nSSP [1, Observation 3.4], and an odd-order double path whose loops all lie on even vertices
does not allow the nSSP [4, Theorem 5.7(Cc)]. We do not claim a classification of which allowable
patterns require the nSSP.

\section{Definitions and notation}


For $n\geq1$, write $[n]=\{1,\ldots,n\}$. The double path with loop assignment
$L\subseteq[n]$ is the directed graph

\[
  P_{n,L}:\qquad
  E(P_{n,L})=
  \{(j,j+1),(j+1,j):1\leq j<n\}
  \cup\{(\ell,\ell):\ell\in L\}.
\]

Let $\mathcal M(P_{n,L})$ be the set of real matrices whose nonzero pattern is exactly
$P_{n,L}$. A pattern \textbf{allows} the nSSP if at least one matrix in its pattern class has the nSSP; it
\textbf{requires} the nSSP if every matrix in the class has it [4].

The positive constructions below are symmetric. Fix $i\in[n]$, a nonzero loop value $a$, and
nonzero edge weights $b_1,\ldots,b_{n-1}$. Put

\[
  A=A_i(a,b)
  =aE_{ii}+\sum_{j=1}^{n-1}b_j(E_{j,j+1}+E_{j+1,j}).       \tag{2.1}
\]

Thus $A\in\mathcal M(P_{n,\{i\}})$. We call such matrices symmetric Jacobi-type matrices; no
positivity assumption on the nonzero $b_j$ is needed.

Let

\[
  J=A-aE_{ii}.
\]

Deleting vertex $i$ from $J$ leaves a left arm of order $i-1$ and a right arm of order
$n-i$. Denote their characteristic polynomials by $L(x)$ and $R(x)$, respectively. An empty arm
has characteristic polynomial $1$.

For a zero-diagonal weighted path, its monic continuants satisfy

\[
  K_0(x)=1,\qquad K_1(x)=x,\qquad
  K_s(x)=xK_{s-1}(x)-c_{s-1}^2K_{s-2}(x).                \tag{2.2}
\]

Consequently $K_s(-x)=(-1)^sK_s(x)$. Such a path has simple spectrum, symmetric about zero; zero is
an eigenvalue exactly when its order is odd.

\section{The nSSP as a spectral-coefficient Jacobian}


We first record the Jacobian bridge in the special form needed here. General relations between strong
properties, transversality, and Jacobian methods are developed in [1--3].

\begin{lemma}[pattern functionals and the spectral Jacobian]
\label{lemma:3-1}

Let $A=A_i(a,b)$ be as in (2.1). Let $\Phi(a,b_1,\ldots,b_{n-1})$ be the vector of the $n$
non-leading coefficients of $\det(xI-A)$. Then the following are equivalent:

\begin{enumerate}
\item $A$ has the nSSP;
\item $D\Phi$ is nonsingular;
\item the $n$ functionals
   \[
     p\longmapsto p(A)_{ii},\qquad
     p\longmapsto p(A)_{j,j+1}\quad(1\leq j<n)
   \]
   are linearly independent on $\mathbb R[x]_{<n}$.
\end{enumerate}
\end{lemma}

\begin{proof}

Every real symmetric irreducible tridiagonal matrix has simple spectrum. Hence its real centralizer is

\[
  \{p(A):p\in\mathbb R[x]_{<n}\}.
\]

In particular, every matrix commuting with $A$ is symmetric. If
$[A,X^{\mathsf T}]=0$, then $X^{\mathsf T}=p(A)$ for a unique polynomial of degree less than
$n$, and symmetry gives $X=p(A)$. Since the nonzero entries of $A$ occur exactly at $(i,i)$
and on the two orientations of each path edge, the equation $A\circ X=0$ is equivalent to

\[
  p(A)_{ii}=0,\qquad p(A)_{j,j+1}=0\quad(1\leq j<n).
\]

Thus (1) and (3) are equivalent.

For $1\leq k\leq n$, define the normalized moments

\[
  m_k(A)=\frac1k\operatorname{tr}(A^k).
\]

Differentiation gives

\[
  \frac{\partial m_k}{\partial a}=(A^{k-1})_{ii},
  \qquad
  \frac{\partial m_k}{\partial b_j}=2(A^{k-1})_{j,j+1}. \tag{3.1}
\]

In the polynomial basis $1,x,\ldots,x^{n-1}$, this moment Jacobian is the transpose of the
functional matrix in (3), with each edge column multiplied by the nonzero scalar $2$. Newton's
identities give an invertible triangular change of coordinates between the first $n$ power sums and
the $n$ non-leading characteristic coefficients. Hence the moment Jacobian and $D\Phi$ have the
same rank, proving the equivalence of (2) and (3). \qedhere

\end{proof}

\section{The spectral-separation criterion}


We now characterize the nSSP for every matrix in the symmetric single-loop family (2.1).

\begin{theorem}[single-loop spectral separation]
\label{theorem:4-1}

Let $A=A_i(a,b)$, where $a\neq0$ and every $b_j\neq0$. Then

\[
  A\text{ has the nSSP}
  \quad\Longleftrightarrow\quad
  \gcd(L,R)=1,
\]

equivalently, the left and right arms obtained by deleting $i$ have disjoint spectra.

\end{theorem}

\begin{proof}[Proof: disjoint arm spectra imply the nSSP]

Let

\[
  P(x)=\det(xI-J),\qquad Q(x)=L(x)R(x).
\]

Expanding at the unique diagonal perturbation gives

\[
  \det(xI-A)=P(x)-aQ(x).                                  \tag{4.1}
\]

Consider an infinitesimal parameter perturbation
$(\delta a,\delta b_1,\ldots,\delta b_{n-1})$ in the kernel of $D\Phi$. Equation (4.1) yields

\[
  0=\delta P-(\delta a)Q-a\,\delta Q.                     \tag{4.2}
\]

By (2.2), $P$ and $\delta P$ contain only powers with the parity of $n$, whereas $Q$ and
$\delta Q$ contain only powers with the parity of $n-1$. The two parity parts of (4.2) therefore
vanish separately:

\[
  \delta P=0,\qquad (\delta a)Q+a\,\delta Q=0.            \tag{4.3}
\]

The polynomial $Q$ is monic of degree $n-1$, while its first variation has degree strictly less
than $n-1$. Comparing leading coefficients in (4.3) gives $\delta a=0$, and then
$\delta Q=0$.

Assume $\gcd(L,R)=1$. Since

\[
  0=\delta Q=(\delta L)R+L(\delta R),                    \tag{4.4}
\]

Euclid's lemma and the strict degree bounds
$\deg\delta L<\deg L$, $\deg\delta R<\deg R$ imply

\[
  \delta L=\delta R=0.                                   \tag{4.5}
\]

The conclusion remains literal when an arm is empty: its polynomial is the constant $1$ and has zero
variation.

Let $L^-$ and $R^-$ be the characteristic polynomials obtained by deleting the arm vertex adjacent
to $i$; for a one-vertex arm the resulting polynomial is $1$. Define, when the corresponding arm
exists,

\[
  U=b_{i-1}^2L^-,
  \qquad
  V=b_i^2R^-.
\]

The absent term is omitted at an endpoint. Expanding the zero-diagonal full path at vertex $i$ gives

\[
  P=xLR-UR-LV.                                            \tag{4.6}
\]

Using $\delta P=\delta L=\delta R=0$ in the first variation of (4.6), we obtain

\[
  (\delta U)R+L(\delta V)=0.                              \tag{4.7}
\]

Again coprimality and
$\deg\delta U<\deg L$, $\deg\delta V<\deg R$ imply that every existing term satisfies

\[
  \delta U=\delta V=0.                                   \tag{4.8}
\]

For example, if the left arm is nonempty, the leading coefficient of
$\delta(b_{i-1}^2L^-)$ is $\delta(b_{i-1}^2)$. Thus (4.8) first gives

\[
  \delta(b_{i-1}^2)=0,\qquad \delta L^-=0.
\]

Apply the continuant recurrence (2.2) successively to the rooted pair $(L,L^-)$. If two consecutive
continuants have zero first variation, comparison of the leading coefficient in

\[
  \delta\bigl(c^2K_{s-2}\bigr)=0
\]

first gives $\delta(c^2)=0$ and then $\delta K_{s-2}=0$. Finite induction recovers every internal
left-arm edge square. The right arm is identical. Hence $\delta(b_j^2)=0$ for all $j$. Since every
$b_j\neq0$, all $\delta b_j=0$; together with $\delta a=0$, the Jacobian kernel is trivial.
Lemma 3.1 now gives the nSSP.

\end{proof}

\begin{proof}[Proof: a common arm eigenvalue obstructs the nSSP]

Suppose the arms share an eigenvalue $\mu$. Both arms are then nonempty. Choose corresponding arm
eigenvectors. Their coordinates at the endpoints adjacent to $i$ are nonzero, because a nonzero
eigenvector of an irreducible tridiagonal matrix cannot vanish at an endpoint. Scale the two vectors so
that their contributions through the two edges incident with $i$ cancel. After embedding them in
$\mathbb R^n$ and putting zero at coordinate $i$, we obtain a nonzero vector $z$ satisfying

\[
  Az=\mu z,\qquad z_i=0.                                 \tag{4.9}
\]

Let $S$ be a diagonal sign matrix alternating along the full path, so
$S_jS_{j+1}=-1$. The zero-diagonal matrix $J$ anticommutes with $S$. Because $z_i=0$, (4.9)
also gives $Jz=\mu z$, and therefore

\[
  A(Sz)=-\mu(Sz).                                        \tag{4.10}
\]

If $\mu\neq0$, set

\[
  X=zz^{\mathsf T}+(Sz)(Sz)^{\mathsf T}.                 \tag{4.11}
\]

Each rank-one term is a spectral projector up to scale, so $[A,X]=0$. At every path edge
$(j,j+1)$,

\[
  X_{j,j+1}=z_jz_{j+1}\bigl(1+S_jS_{j+1}\bigr)=0,
\]

and $X_{ii}=0$. Thus $A\circ X=0$, while $X\neq0$, so $A$ fails the nSSP.

If $\mu=0$, each arm must have odd order. Its zero mode is supported on alternating vertices, hence
the product of its two coordinates across every arm edge is zero. The vector $z$ in (4.9) inherits
this property, and

\[
  X=zz^{\mathsf T}
\]

already satisfies $[A,X]=0$, $A\circ X=0$, and $X\neq0$. This proves the converse. \qedhere

\end{proof}

\begin{remark}
\label{remark:4-2}

The proof of the sufficient direction is a local, first-order recovery argument. It does not infer
Jacobian nonsingularity merely from global uniqueness in the three-spectra inverse problem. Conversely,
the explicit matrices (4.11) and $zz^{\mathsf T}$ show directly why a shared arm eigenvalue violates
the nSSP equations.

\end{remark}

\section{Single-loop existence}


The structural criterion turns existence into a parity question plus a finite collision-avoidance step.

\begin{lemma}[rational spectral separation]
\label{lemma:5-1}

Suppose the two arm orders have opposite parity, or both are even. Then the edge weights in (2.1) can be
chosen nonzero and rational so that $\gcd(L,R)=1$.

\end{lemma}

\begin{proof}

A zero-diagonal irreducible path of odd order has one simple zero eigenvalue, while one of even order has
none. Under the stated parity hypothesis, the two arms do not both contain zero.

Choose arbitrary nonzero rational internal edge weights on both arms. If necessary, multiply every
internal edge weight of the right arm by a positive rational number $c$. Every nonzero right-arm
eigenvalue is then multiplied by $c$. A collision with a nonzero left-arm eigenvalue can occur only for
one of finitely many positive ratios $|\lambda/\nu|$, where $\lambda$ and $\nu$ range over the
nonzero spectra of the two unscaled arms. Choose $c\in\mathbb Q_{>0}$ outside this finite set. Empty
arms and one-vertex arms require no scaling. Assign arbitrary nonzero rational weights to the one or two
edges, when present, that join the arms to vertex $i$. These weights do not enter $L$ or $R$. The
two arm spectra are then disjoint. \qedhere

\end{proof}

\begin{theorem}[single-loop classification]
\label{theorem:5-2}

For every $n\geq1$ and $i\in[n]$,

\[
  P_{n,\{i\}}\text{ allows the nSSP}
  \quad\Longleftrightarrow\quad
  n\text{ is even or }i\text{ is odd}.                   \tag{5.1}
\]

\end{theorem}

\begin{proof}

The arm orders are $i-1$ and $n-i$.

\begin{itemize}
\item If $n$ is even, these orders have opposite parity.
\item If $n$ is odd and $i$ is odd, both orders are even.
\item If $n$ is odd and $i$ is even, both orders are odd.
\end{itemize}
In the first two cases, Lemma 5.1 gives rational arm weights with disjoint spectra. Choose the loop and
the one or two edges incident with $i$ to be any nonzero rational numbers. Theorem 4.1 produces a
symmetric matrix in $\mathcal M(P_{n,\{i\}})$ with the nSSP.

In the third case, nonexistence is stronger than the failure of our symmetric construction:
Koljan\v{c}i\'{c}--Oblak [4, Theorem 5.7(Cc)] proves that for odd $n$, every double path whose loops lie only
on even vertices does not allow the nSSP. Applying that theorem to $L=\{i\}$ gives the reverse
implication. The case $n=1$ is included directly: the one-by-one matrix $[a]$, $a\neq0$, has the
nSSP. \qedhere

\end{proof}

\section{All loop assignments}


We now combine the single-loop construction with the Superpattern Lemma and the two known negative
cases.

\begin{theorem}[allow/not-allow classification]
\label{theorem:6-1}

Let $n\geq1$ and $L\subseteq[n]$. Then

\[
  P_{n,L}\text{ allows the nSSP}
  \quad\Longleftrightarrow\quad
  L\neq\varnothing
  \quad\text{and}\quad
  \bigl(n\text{ is even or }L\cap\{1,3,5,\ldots\}\neq\varnothing\bigr). \tag{6.1}
\]

\end{theorem}

\begin{proof}

Assume first that the condition on the right holds. If $n$ is even, choose any $i\in L$; if
$n$ is odd, choose an odd $i\in L$. By Theorem 5.2, the single-loop pattern
$P_{n,\{i\}}$ allows the nSSP. Since $P_{n,L}$ is a superpattern, the Superpattern Lemma
[3, 4] shows that $P_{n,L}$ also allows the nSSP.

Conversely, if $L=\varnothing$, a loopless pattern does not allow the nSSP
[1, Observation 3.4]. In the remaining case, if $n\geq3$ is odd and $L$ contains no odd vertex, then
every loop lies on an even vertex, and [4, Theorem 5.7(Cc)] again says that $P_{n,L}$ does not allow the
nSSP. These are exactly the failures of the right-hand condition. \qedhere

\end{proof}

\begin{corollary}[Question 5.11]
\label{corollary:6-2}

If $n$ is even and $L\neq\varnothing$, then $P_{n,L}$ allows the nSSP.

\end{corollary}

\begin{remark}[scope of the classification]
\label{remark:6-3}

Theorem 6.1 gives a complete \textbf{allow/not-allow} classification. It does not decide whether an allowable
pattern requires the nSSP or merely contains at least one matrix having it. Thus it answers Question 5.11
but does not complete the three-way classification requested in [4, Problem 1.6].

\end{remark}

\section{Relation to previous work}


The novelty boundary is deliberately narrow.

\begin{enumerate}
\item The nSSP, the Superpattern Lemma, and the similarity-transversality viewpoint are prior theory
   [1, 3]. General comparisons between strong properties and Jacobian methods are also available [2].
   Lemma 3.1 is included as a self-contained specialization, not claimed as a new general framework.
\item Continuant recurrences, simple spectrum of irreducible symmetric tridiagonal matrices, and parity of
   zero-diagonal path spectra are classical.
\item Michor--Teschl [5] already identify disjoint spectra of the two deleted arms as the condition under
   which three spectra reconstruct a Jacobi matrix without additional mixing information. We do not
   claim spectral separation or three-spectra recovery as new.
\item Koljan\v{c}i\'{c}--Oblak [4] provide the double-path formulation and Question 5.11, as well as general
   positive and negative subfamilies. In particular, their Theorem 5.5 covers all initial-segment loop
   assignments and hence all endpoint single loops; Theorem 5.7(B) covers $P_{n,\{2\}}$ for even
   $n$; Theorem 5.7(Cc) gives the odd-order negative family used here; and Examples 5.9--5.10 give
   further internal single-loop cases.
\end{enumerate}
To the best of our knowledge, the items not directly covered by these sources are Theorem 4.1 as an nSSP
criterion with an explicit common-spectrum obstruction, the completion in Theorem 5.2 of the single-loop
allow/not-allow classification beyond the subfamilies just listed, and the combined classification in
Theorem 6.1. Thus the contribution is a unified spectral criterion and a completion of known positive and
negative subfamilies, not the first positive result for single-loop double paths. This statement records
the result of a targeted public-literature search; it is not an assertion that all unpublished or
unindexed work has been excluded.

\section{Computational calibration and reproducibility}


The proof above is general and does not depend on finite enumeration. Three complementary scripts are
retained to expose indexing, parity, boundary, and implementation errors. The first and second reuse the
weighted-path construction and modular-rank routines from the third, so they are not three
implementation-independent verifiers; they test different identities and controls over a shared core:

\begin{enumerate}
\item \texttt{verify\_nssp\_spectral\_separation.py} compares polynomial coprimality, mod-$p$ rank calibration, and
   the parity classification at all 119 single-loop positions with $2\leq n\leq15$.
\item \texttt{audit\_nssp\_spectral\_jacobian.py} compares the characteristic-coefficient Jacobian, moment Jacobian,
   and full verification matrix on 77 weighted samples. It also constructs exact commuting witnesses for
   one nonzero common eigenvalue and one common zero eigenvalue.
\item \texttt{verify\_nssp\_weighted\_paths.py --max-n 32} checks all 272 naturally labelled single-loop positions
   at even orders through $32$ for the discovery formula $b_j=j$.
\end{enumerate}
The third computation is only a finite check of one convenient weight formula. The general existence
proof uses Lemma 5.1 and makes no claim that $b_j=j$ works for every even order.

From the manuscript entry directory, the checks require Python 3.11 or newer and only the standard
library:

\begin{Verbatim}[fontsize=\small,breaklines=true,breakanywhere=true]
python -X utf8 verification/verify_nssp_weighted_paths.py --max-n 32
python -X utf8 verification/verify_nssp_spectral_separation.py
python -X utf8 verification/audit_nssp_spectral_jacobian.py
\end{Verbatim}

A full-rank result modulo a prime certifies full rank over $\mathbb Q$ and $\mathbb R$. A rank
deficiency modulo a prime is used only as calibration and is not treated as a real-rank proof; the
negative direction is instead certified analytically by Theorem 4.1.

\section{Discussion}


The classification illustrates a useful division of labour between pattern theory and inverse spectral
structure. Pattern results supply the Superpattern Lemma and the sharp nonexistence region. Inside the
remaining region, a symmetric Jacobi-type witness is enough, and its nSSP reduces to separation of two
small spectra. The unavoidable zero mode explains the entire parity obstruction; every nonzero collision
can be removed by a one-parameter rational scaling.

Two questions remain outside the present result. First, one can ask which of the allowable loop
assignments require the nSSP. That problem concerns all matrices in a generally non-symmetric pattern and
cannot be settled by the symmetric existence witnesses used here. Second, it may be useful to determine
whether the spectral-separation criterion extends to broader tree-like patterns with one loop, where
deleting the looped vertex produces more than two branches and common branch spectra may interact in a
higher-dimensional way.

\section*{Declaration of AI-assisted research and writing}


This project used OpenAI Codex extensively for problem exploration, symbolic derivation, adversarial
proof checking, literature-search assistance, verification-code generation, and manuscript drafting.
The retained claims are presented through human-readable arguments and reproducible checks, but AI output
is not treated as an authority and external expert review remains pending. Codex is a tool, not an author.
The named author, Carptopus, accepts responsibility for the manuscript's accuracy, originality,
citations, and integrity.

\section*{License}


Copyright 2026 Carptopus. This manuscript is licensed under the
\href{https://creativecommons.org/licenses/by/4.0/}{Creative Commons Attribution 4.0 International License (CC BY 4.0)}.
Reusers must provide appropriate attribution and indicate whether changes were made. No warranties are
given.

\begin{thebibliography}{9}


\bibitem{ref1} M. Arav, F. J. Hall, H. van der Holst, Z. Li, A. Mathivanan, J. Pan, H. Xu and Z. Yang,
   \emph{Advances on similarity via transversal intersection of manifolds}, Linear Algebra and its
   Applications 721 (2025), 102--121, \url{https://doi.org/10.1016/j.laa.2024.05.013}.
\bibitem{ref2} M. Catral, S. Fallat, H. Gupta and J. C.-H. Lin,
   \emph{The Strong Spectral Property and the Jacobian Method for Weighted Laplacian Matrices},
   arXiv:2602.18999, 2026, \url{https://arxiv.org/abs/2602.18999}.
\bibitem{ref3} S. M. Fallat, H. T. Hall, J. C.-H. Lin and B. L. Shader,
   \emph{The bifurcation lemma for strong properties in the inverse eigenvalue problem of a graph},
   Linear Algebra and its Applications 648 (2022), 70--87.
\bibitem{ref4} S. Koljan\v{c}i\'{c} and P. Oblak,
   \emph{Combinatorial aspects of the non-symmetric strong spectral property for graphs},
   arXiv:2604.21552v1, 2026, \url{https://arxiv.org/abs/2604.21552}.
\bibitem{ref5} J. Michor and G. Teschl,
   \emph{Reconstructing Jacobi Matrices from Three Spectra}, in J. Janas, P. Kurasov and S. Naboko (eds.),
   \emph{Spectral Methods for Operators of Mathematical Physics}, Operator Theory: Advances and Applications
   154, Birkhäuser, Basel, 2004, 151--154,
   \url{https://doi.org/10.1007/978-3-0348-7947-7_9}.
\end{thebibliography}

\end{document}
